\subsection{Report Structure}

\noindent This report is organized as follows:

\begin{itemize}
    \item \textbf{Chapter 1 -- Introduction:} Introduces the motivation behind aBridgeAI, identifying the limitations of traditional educational systems in delivering personalised, competency-driven learning at institutional scale. Defines the project goals, delimits the system scope, and articulates the academic and practical significance of the proposed solution.
    \item \textbf{Chapter 2 -- Related Works:} Reviews existing educational technologies and mock interview methodologies, identifying their current limitations and establishing the theoretical need for an integrated, context-aware solution.
    \item \textbf{Chapter 3 -- Theoretical Background:} Establishes the theoretical foundation for the core technologies employed in this project, including LLMs, RAG pipelines, Vector Databases, and modern web development frameworks.
    \item \textbf{Chapter 4 -- System Analysis and Requirements:} Details the systematic extraction of software requirements. It encompasses the definition of user roles, comprehensive Functional Requirements (FRs), quantifiable Non-Functional Requirements (NFRs), and Unified Modeling Language (UML) artifacts such as Use Case and Activity diagrams.
    \item \textbf{Chapter 5 -- System Design:} Presents the high-level architectural blueprint of the system. It covers the system topology, PostgreSQL with pgvector, Neo4j, Redis, Garage object storage, LiveKit interview delivery, and User Interface/User Experience (UI/UX) wireframes.
    \item \textbf{Chapter 6 -- Implementation and Testing:} Discusses the practical realization of the proposed pipeline, including the native LiveKit interview runtime and its automated-test strategy.
    \item \textbf{Chapter 7 -- Evaluation:} Presents the measured frontend, backend, and quiz-generation evaluation evidence, together with FRS inventory and test-case traceability.
    \item \textbf{Chapter 8 -- Conclusion:} Summarizes the technical contributions and educational impacts of the project, acknowledges current system limitations, and outlines strategic directions for future enhancements.
\end{itemize}